\Article{19}
\textsc{Satz.} \textit{Bleibt alles wie oben, so lässt sich um den Mittelpunkt $C$ ein Kreis beschreiben,}
\SourceLeaf{030}
\textit{auf dessen Umfang $2m$ Punkte mit $T=0$ und ebenso viele mit $U=0$ liegen, und zwar so, dass jeder der letzteren zwischen zweien der ersteren liegt.}

Die Summe aller positiv genommenen Koeffizienten $A,B$ usw. $K,L,M$ sei $=S$, und man wähle $R$ zugleich $>S\BareRoot{2}$ und $>1$\HistoricalNote{9}. Dann behaupte ich, dass für den mit dem Radius $R$ beschriebenen Kreis notwendig gilt, was der Satz ausspricht. Der Kürze halber bezeichne man den Punkt dieses Umfangs, der $\frac1m45$ Grad von dessen Schnitt mit dem linken Teil der Achse entfernt ist, oder für den $\varphi=\frac1m45^\circ$ gilt, mit (1); ebenso den Punkt im Abstand $\frac3m45^\circ$ von diesem Schnitt, oder für den $\varphi=\frac3m45^\circ$ gilt, mit (3); ferner den Punkt mit $\varphi=\frac5m45^\circ$ mit (5) usw. bis zu $(8m-1)$. Dieser liegt $\frac{8m-1}m45$ Grad von jenem Schnitt entfernt, wenn man stets in derselben Richtung fortschreitet (oder $\frac1m45^\circ$ in der entgegengesetzten Richtung). So erhält man insgesamt $4m$ Punkte auf dem Umfang in gleichen Abständen. Zwischen $(8m-1)$ und (1) liegt ein Punkt mit $T=0$; ebenso liegt je ein solcher Punkt zwischen (3) und (5), zwischen (7) und (9), zwischen (11) und (13) usw., sodass ihre Anzahl $2m$ beträgt. Auf dieselbe Weise liegen je einzelne Punkte mit $U=0$ zwischen (1) und (3), zwischen (5) und (7), zwischen (9) und (11); ihre Anzahl ist also ebenfalls $=2m$. Außer diesen $4m$ Punkten gibt es schließlich auf dem ganzen Umfang keine weiteren, in denen entweder $T$ oder $U$ den Wert $=0$ hat.

\textit{Beweis.} I. Im Punkt (1) ist $m\varphi=45^\circ$ und daher
\[ T=R^{m-1}(R\BareRoot{\tfrac12}+A\sin(m-1)\varphi+\tfrac BR\sin(m-2)\varphi+\text{usw.}+\tfrac L{R^{m-2}}\sin\varphi)
\]
Die Summe $A\sin(m-1)\varphi+\frac BR\sin(m-2)\varphi$ usw. kann aber gewiss nicht größer als $S$ sein und ist daher notwendig kleiner als $R\BareRoot{\frac12}$. Daraus folgt, dass der Wert von $T$ in diesem Punkt gewiss positiv ist. Umso mehr hat also $T$ einen positiven Wert, wenn $m\varphi$ zwischen $45^\circ$ und $135^\circ$ liegt, d.\,h. der Wert von $T$ ist vom Punkt (1) bis (3) stets positiv. Aus demselben Grund hat $T$ vom Punkt (9) bis (11) überall einen positiven Wert und allgemein von jedem Punkt $(8k+1)$ bis $(8k+3)$, wobei $k$ eine beliebige ganze Zahl bezeichnet. Ebenso hat $T$ überall zwischen (5) und (7), zwischen (13) und (15) usw. und allgemein zwischen $(8k+5)$ und $(8k+7)$ einen negativen Wert; in allen diesen Intervallen kann es also nirgends $=0$ sein. Da dieser Wert aber in (3) positiv und in (5) negativ ist, muss er irgendwo zwischen (3) und (5) $=0$ sein; ebenso irgendwo zwischen (7) und (9), zwischen (11) und (13) usw.
\SourceLeaf{031}
bis einschließlich des Intervalls zwischen $(8m-1)$ und (1), sodass insgesamt in $2m$ Punkten $T=0$ gilt. Q. E. P.

II. Dass es außer diesen $2m$ Punkten keine weiteren mit dieser Eigenschaft gibt, erkennt man folgendermaßen. Da zwischen (1) und (3), zwischen (5) und (7) usw. keine solchen Punkte liegen, könnte es mehr davon nur geben, wenn in irgendeinem Intervall zwischen (3) und (5) oder zwischen (7) und (9) usw. mindestens zwei lägen. Dann müsste aber $T$ irgendwo in demselben Intervall ein \textit{Maximum} oder \textit{Minimum} sein, und daher wäre $\frac{\mathrm dT}{\mathrm d\varphi}=0$. Es ist jedoch $\frac{\mathrm dT}{\mathrm d\varphi}=mR^{m-2}(R\cos m\varphi+\frac{m-1}mA\cos(m-1)\varphi+\text{usw.})$, und $\cos m\varphi$ ist zwischen (3) und (5) stets negativ und $>\BareRoot{\frac12}$. Daraus erkennt man leicht, dass $\frac{\mathrm dT}{\mathrm d\varphi}$ in diesem ganzen Intervall eine negative Größe ist; ebenso ist es zwischen (7) und (9) überall positiv, zwischen (11) und (13) negativ usw. Es kann also in keinem dieser Intervalle 0 sein, und die Annahme kann nicht bestehen. Folglich usw. Q. E. S.

III. Ganz ebenso beweist man, dass $U$ überall zwischen (3) und (5), zwischen (11) und (13) usw. und allgemein zwischen $(8k+3)$ und $(8k+5)$ einen negativen Wert hat, dagegen zwischen (7) und (9), zwischen (15) und (17) usw. und allgemein zwischen $(8k+7)$ und $(8k+9)$ einen positiven. Daraus folgt sofort, dass irgendwo zwischen (1) und (3), zwischen (5) und (7) usw., also in $2m$ Punkten, $U=0$ gelten muss. In keinem dieser Intervalle kann aber $\frac{\mathrm dU}{\mathrm d\varphi}=0$ sein (was man leicht ebenso wie oben beweist). Daher gibt es auf dem Kreisumfang nicht mehr als jene $2m$ Punkte, in denen $U=0$ gilt. Q. E. T. und Q.

Der Teil des Satzes, wonach es nicht mehr als $2m$ Punkte mit $T=0$ und nicht mehr als $2m$ Punkte mit $U=0$ gibt, lässt sich übrigens auch daraus beweisen, dass die Gleichungen $T=0$, $U=0$ Kurven vom Grad $m$ darstellen. Aus der höheren Geometrie ist bekannt, dass solche Kurven von einem Kreis als Kurve zweiten Grades nicht in mehr als $2m$ Punkten geschnitten werden können.

\Article{20} Beschreibt man um denselben Mittelpunkt einen anderen Kreis mit einem größeren Radius als $R$ und teilt ihn auf dieselbe Weise, so liegt auch auf ihm zwischen den Punkten (3) und (5) ein Punkt mit $T=0$, ebenso zwischen (7) und (9) usw. Man erkennt leicht, dass solche zwischen (3) und (5) auf den beiden Umfängen liegenden Punkte einander umso näher sein müssen, je weniger sich der Radius dieses Kreises vom Radius $R$ unterscheidet. Dasselbe gilt, wenn man einen Kreis mit einem etwas kleineren Radius als $R$ beschreibt, der jedoch größer als $S\BareRoot{2}$ und 1 bleibt. Daraus versteht man ohne Mühe, dass der Umfang des mit dem Radius $R$ beschriebenen Kreises
\SourceLeaf{032}
in dem Punkt zwischen (3) und (5), in dem $T=0$ gilt, tatsächlich von irgendeinem Zweig der ersten Linie \textit{geschnitten} wird; dasselbe gilt für die übrigen Punkte mit $T=0$. Ebenso ist klar, dass der Umfang dieses Kreises in allen $2m$ Punkten mit $U=0$ von irgendeinem Zweig der zweiten Linie geschnitten wird. Diese Folgerungen lassen sich auch so ausdrücken: Beschreibt man um den Mittelpunkt $C$ einen Kreis von der erforderlichen Größe, so treten $2m$ Zweige der ersten Linie und ebenso viele der zweiten in ihn ein, und zwar so, dass je zwei benachbarte Zweige der ersten Linie durch irgendeinen Zweig der zweiten voneinander getrennt werden. Siehe Fig. 2, in der der Kreis nun eine endliche statt einer unendlichen Größe hat. Die den einzelnen Zweigen beigeschriebenen Zahlen dürfen nicht mit den Zahlen verwechselt werden, mit denen ich im vorigen und in diesem Artikel der Kürze halber bestimmte Grenzen auf dem Umfang bezeichnet habe.

\Article{21} Aus dieser gegenseitigen Lage der in den Kreis eintretenden Zweige lässt sich die Notwendigkeit eines Schnittes irgendeines Zweiges der ersten Linie mit einem Zweig der zweiten innerhalb des Kreises auf so viele verschiedene Arten herleiten, dass ich beinahe nicht weiß, welche Methode vor den übrigen zu wählen ist. Am deutlichsten scheint mir die folgende. Bezeichnen wir (Fig. 2) den Punkt des Kreisumfangs, in dem er vom linken Teil der Achse geschnitten wird (der selbst einer der $2m$ Zweige der ersten Linie ist), mit 0; den nächsten Punkt, in dem ein Zweig der zweiten Linie eintritt, mit 1; den darauf folgenden, in dem der zweite Zweig der ersten Linie eintritt, mit 2; und so fort bis $4m-1$. Damit tritt in jedem mit einer geraden Zahl bezeichneten Punkt ein Zweig der zweiten Linie in den Kreis ein, dagegen ein Zweig der zweiten Linie in allen mit ungeraden Zahlen bezeichneten Punkten. Aus der höheren Geometrie ist nun bekannt, dass jede algebraische Kurve (oder jeder einzelne Teil einer algebraischen Kurve, falls sie aus mehreren zusammengesetzt ist) entweder in sich zurückkehrt oder nach beiden Seiten ins Unendliche verläuft. Tritt also irgendein Zweig einer algebraischen Kurve in ein begrenztes Gebiet ein, so muss er es notwendig irgendwo wieder verlassen\HistoricalNote{10}. Daraus schließt man leicht, dass jeder mit einer geraden Zahl\SourceLeaf{033}
 bezeichnete Punkt (oder der Kürze halber \textit{jeder gerade Punkt}) innerhalb des Kreises durch einen Zweig der ersten Linie mit einem anderen geraden Punkt verbunden sein muss; ebenso jeder mit einer ungeraden Zahl bezeichnete Punkt durch einen Zweig der zweiten Linie mit einem anderen solchen Punkt. Obwohl diese Verbindung von je zwei Punkten je nach der Beschaffenheit des Polynoms $X$ sehr verschieden sein und daher nicht allgemein bestimmt werden kann, lässt sich doch leicht beweisen, dass \textit{sie, wie sie auch beschaffen sein mag, stets einen Schnitt der ersten mit der zweiten Linie hervorbringt.}

\Article{22} Der Beweis dieser Notwendigkeit lässt sich wohl am bequemsten durch Widerspruch führen. Nehmen wir also an, die Verbindung je zweier gerader und je zweier ungerader Punkte könne so eingerichtet werden, dass daraus kein Schnitt eines Zweiges der ersten Linie mit einem Zweig der zweiten entsteht. Da die Achse ein Teil der ersten Linie ist, ist offenbar der Punkt 0 mit dem Punkt $2m$ verbunden. Der Punkt 1 kann daher mit keinem jenseits der Achse liegenden Punkt verbunden sein, d.\,h. mit keinem Punkt, der durch eine größere Zahl als $2m$ bezeichnet wird; sonst würde die Verbindungslinie notwendig die Achse schneiden. Nimmt man also an, \mbox{1 sei mit dem Punkt $n$ verbunden, so ist $n<2m$}. Aus einem ähnlichen Grund gilt $n'<n$, wenn man annimmt, 2 sei mit $n'$ verbunden; sonst würde nämlich der Zweig $2\ldots n'$ notwendig den Zweig $1\ldots n$ schneiden. Aus demselben Grund ist der Punkt 3 mit einem der zwischen 4 und $n'$ liegenden Punkte verbunden. Nimmt man an, 3, 4, 5 usw. seien mit $n''$, $n'''$, $n''''$ usw. verbunden, so liegt offenbar $n'''$ zwischen 5 und $n''$, $n''''$ zwischen 6 und $n'''$ usw. Es ist also klar, dass man schließlich zu irgendeinem Punkt $h$ gelangt, der mit dem Punkt $h+2$ verbunden ist. Dann schneidet der Zweig, der im Punkt $h+1$ in den Kreis eintritt, notwendig den Zweig, der die Punkte $h$ und $h+2$ verbindet. Da aber der eine dieser beiden Zweige zur ersten Linie und der andere zur zweiten gehört, ist nun klar, dass die Annahme widersprüchlich ist und daher notwendig irgendwo ein Schnitt der ersten mit der zweiten Linie stattfindet.
\SourceLeaf{034}
Verbindet man dies mit dem Vorhergehenden, so ergibt sich aus allen dargelegten Untersuchungen, dass der Satz, \textit{jedes Polynom in einer Variablen lasse sich in reelle Faktoren ersten oder zweiten Grades zerlegen}, mit voller Strenge bewiesen ist.

\Article{23} Aus denselben Grundsätzen lässt sich übrigens ohne große Schwierigkeit herleiten, dass es nicht nur einen, sondern mindestens $m$ Schnitte der ersten mit der zweiten Linie gibt. Allerdings kann die erste Linie auch von mehreren Zweigen der zweiten im selben Punkt geschnitten werden; in diesem Fall besitzt das Polynom $X$ mehrere gleiche Faktoren. Da es hier jedoch genügt, die Notwendigkeit eines einzigen Schnittes bewiesen zu haben, halte ich mich der Kürze halber nicht länger dabei auf. Aus demselben Grund gehe ich hier auch auf andere Eigenschaften dieser Linien nicht näher ein, etwa darauf, dass die Schnitte stets unter rechten Winkeln erfolgen; oder dass, wenn mehrere Arme beider Kurven in demselben Punkt zusammenkommen, ebenso viele Arme der ersten wie der zweiten Linie vorhanden sind, die abwechselnd liegen und sich unter gleichen Winkeln schneiden usw.

Schließlich bemerke ich, dass es keineswegs unmöglich ist, den vorigen Beweis, den ich hier auf geometrische Grundsätze gestützt habe, auch in rein analytischer Form darzustellen. Doch glaubte ich, die hier entwickelte Darstellung werde weniger abstrakt ausfallen und die eigentliche Kraft des Beweises werde hier viel deutlicher vor Augen geführt, als von einem analytischen Beweis zu erwarten wäre.

Zum Abschluss werde ich noch eine andere Methode zum Beweis unseres Satzes andeuten. Auf den ersten Blick wird sie nicht nur von dem vorigen Beweis, sondern auch von allen übrigen oben angeführten Beweisen sehr verschieden erscheinen; ihrem Wesen nach ist sie dennoch eigentlich dieselbe wie die von \textsc{d'Alembert}. Den Vergleich mit dieser und die Untersuchung der Parallele zwischen beiden überlasse ich den Sachkundigen, allein für die sie hinzugefügt ist.

\Article{24} Über der Ebene von Fig. 4 denke ich mir die erste und die zweite Fläche in Bezug auf die Achse $CG$ und den festen Punkt $C$ auf dieselbe Weise wie oben beschrieben. Man nehme irgendeinen Punkt auf irgendeinem Zweig der ersten Linie, also einen Punkt mit $T=0$ (etwa einen beliebigen Punkt $M$ auf der Achse). Gilt in ihm nicht auch $U=0$, so gehe man von diesem Punkt auf der ersten Linie in der Richtung weiter, in der der absolute\SourceLeaf{035}
 Betrag von $U$ abnimmt. Nimmt im Punkt $M$ der absolute Betrag von $U$ nach beiden Seiten ab, so ist es gleichgültig, in welcher Richtung man weitergeht; was zu tun ist, wenn $U$ nach beiden Seiten wächst, werde ich sogleich zeigen. Es ist also klar, dass man, wenn man stets auf der ersten Linie weitergeht, notwendig schließlich zu einem Punkt gelangt, in dem $U=0$ gilt, oder zu einem solchen, in dem der Wert von $U$ ein Minimum wird, etwa zum Punkt $N$. Im ersten Fall ist das Gesuchte gefunden. Im zweiten lässt sich beweisen, dass sich in diesem Punkt mehrere Zweige der ersten Linie schneiden (und zwar eine gerade Anzahl von Zweigen), von denen die Hälfte so beschaffen ist, dass beim Abbiegen in einen von ihnen (nach dieser oder jener Seite) der Wert von $U$ weiterhin abnimmt. (Den Beweis dieses Satzes, der eher lang als schwierig ist, muss ich der Kürze halber weglassen.) Auf diesem Zweig kann man also wieder weitergehen, bis $U$ entweder $=0$ wird (wie es in Fig. 4 in $P$ geschieht) oder wiederum ein Minimum. Dann biegt man erneut ab und gelangt notwendig schließlich zu einem Punkt, in dem $U=0$ gilt.

Gegen diesen Beweis könnte man den Zweifel vorbringen, ob es nicht möglich sei, dass, wie weit man auch fortschreitet und obwohl der Wert von $U$ stets abnimmt, diese Abnahmen doch immer langsamer werden und jener Wert irgendeine Grenze nirgends erreicht. Dieser Einwand entspräche dem vierten in Artikel 6. Es wäre aber nicht schwer, eine Grenze anzugeben, jenseits deren sich der Wert von $U$ notwendig nicht nur immer schneller ändern muss, sondern auch nicht mehr \textit{abnehmen} kann. Der Wert 0 muss also notwendig schon vor dem Erreichen dieser Grenze vorgekommen sein. Dies und das Übrige, was ich in diesem Beweis nur andeuten konnte, behalte ich mir für eine andere Gelegenheit zu einer ausführlicheren Behandlung vor.

\begin{center}
\rule{8cm}{0.3pt}\\[8pt]
\textit{Die Grundsätze, auf denen dieser Beweis beruht, haben wir Anfang Oktober} 1797 \textit{entdeckt.}\\[8pt]
\rule{5cm}{0.3pt}
\end{center}
